\begin{theorem}[输出计数的准幂极限定律与显式一阶 Edgeworth 系数]\label{thm:xi-real-input-40-output-edgeworth-quasipowers}
沿用推论 \ref{cor:real-input-40-ground-entropy} 与命题 \ref{prop:real-input-40-output-cumulants-closed} 的记号。令
\[
Z_n(u)=\mathbf{1}^{\mathsf T}M(u)^n\mathbf{1}
=\sum_{k=0}^n A_{n,k}u^k,
\qquad
\mathbb P(X_n=k)=\frac{A_{n,k}}{Z_n(1)}.
\]
再记
\[
\mu:=P'(0),
\qquad
\kappa_r:=P^{(r)}(0),
\qquad
\widetilde X_n:=\frac{X_n-n\mu}{\sqrt{n\kappa_2}}.
\]
则 \(\widetilde X_n\Longrightarrow \mathcal N(0,1)\)。并且在中心窗口内一致成立
\[
\mathbb P(\widetilde X_n\le x)
=
\Phi(x)
+\frac{\kappa_3}{6\kappa_2^{3/2}\sqrt n}(1-x^2)\phi(x)
+O(n^{-1}).
\]
特别地，
\[
\frac{\kappa_3}{6\kappa_2^{3/2}}
\approx -0.2460802740027555.
\]
此外，标准化偏度与超额峰度满足
\[
\mathrm{Skew}(X_n)
=
\frac{\kappa_3}{\kappa_2^{3/2}}\frac{1}{\sqrt n}
+O(n^{-1})
=
-\frac{1.4764816440165331}{\sqrt n}+O(n^{-1}),
\]
\[
\mathrm{ExKurt}(X_n)
=
\frac{\kappa_4}{\kappa_2^2}\frac{1}{n}
+O(n^{-3/2})
=
\frac{2.3928814165461985}{n}+O(n^{-3/2}).
\]
因此第一 Edgeworth 修正的符号与尺度都被单一代数常数 \(\kappa_3/(6\kappa_2^{3/2})\) 完全锁定；而离散分位点的精细描述仍需再附加独立的 lattice 校正。
\end{theorem}
\begin{proof}
由主特征值在 \(u=1\) 处的单性与解析扰动，存在 \(\theta=0\) 的复邻域、解析且不为零的函数 \(C(\theta)\) 与常数 \(0<\beta<1\)，使得
\[
Z_n(e^\theta)
=
C(\theta)e^{nP(\theta)}\bigl(1+O(\beta^n)\bigr)
\]
在该邻域内一致成立。于是
\[
\log \mathbb E[e^{\theta X_n}]
=
\log Z_n(e^\theta)-\log Z_n(1)
=
n\bigl(P(\theta)-P(0)\bigr)
+\log\frac{C(\theta)}{C(0)}
+O(\beta^n).
\]
这正是附录 \ref{app:pressure-analytic} 中有限状态解析加权矩阵的标准准幂情形，故中心极限定理与一阶 Edgeworth 展开在中心窗口内成立；其首项系数恰为
\(
\kappa_3/(6\kappa_2^{3/2})
\)。

另一方面，由上式在 \(\theta=0\) 处逐阶求导可得
\[
\operatorname{Var}(X_n)=n\kappa_2+O(1),
\qquad
\kappa_3(X_n)=n\kappa_3+O(1),
\qquad
\kappa_4(X_n)=n\kappa_4+O(1),
\]
从而
\[
\mathrm{Skew}(X_n)
=
\frac{\kappa_3(X_n)}{\operatorname{Var}(X_n)^{3/2}}
=
\frac{\kappa_3}{\kappa_2^{3/2}}\frac{1}{\sqrt n}
+O(n^{-1}),
\]
\[
\mathrm{ExKurt}(X_n)
=
\frac{\kappa_4(X_n)}{\operatorname{Var}(X_n)^2}
=
\frac{\kappa_4}{\kappa_2^2}\frac{1}{n}
+O(n^{-3/2}).
\]
最后，把命题 \ref{prop:real-input-40-output-cumulants-closed} 的闭式代入，并用定理 \ref{thm:xi-real-input-40-output-edgeworth-nongaussian} 中已给出的数值化简，即得各个常数。证毕。
\end{proof}
